\begin{center}
    {\LARGE \textbf{2006物理化学}} \\[2ex]
    {\large 时量180分钟 \quad 满分90分}
\end{center}

\vspace{0.5cm}
\noindent\textbf{注意事项:}
\begin{enumerate}[label=\arabic*., parsep=0pt, itemsep=0pt]
    \item 答题前,考生先将自己的姓名、学号、考场号、座位号填写在答题卡上,将条形码准确粘贴在考生信息条形码粘贴区。
    \item 考试结束后,将本试卷与答题纸一并收回。
    \item 回答各个问题时,将答案写在答题纸上,写在本试卷上无效。
\end{enumerate}

\vspace{0.5cm}
\noindent\textbf{1. (15 pts.)}

$1\ \mathrm{mol}$单原子分子理想气体,始态为$202650\ \mathrm{Pa},11.2\ \mathrm{dm}^3$,经$pT=$常数的可逆过程压缩到终态为$405300\ \mathrm{Pa}$,求
\begin{enumerate}[label=(\arabic*), itemsep=1ex]
    \item 终态的体积和温度；
    \item $\Delta U$和$\Delta H$；
    \item 所做的功。
\end{enumerate}

\vspace{0.5cm}
\noindent\textbf{2. (15 pts.)}

$\ce{Na}$在汞齐中的活度$a_2$符合$\ln a_2 = \ln x_2 + 35.7x_2$,$x_2$为汞齐中$\ce{Na}$的摩尔分数。求$x_2=0.04$时,汞齐中$\ce{Hg}$之活度$a_1$(汞齐中只有$\ce{Na}$及$\ce{Hg}$)。

\vspace{0.5cm}
\noindent\textbf{3. (15 pts.)}

$\ce{CO_{2}}$的固态和液态的蒸气压分别由以下两个方程给出:
\[
\lg(p_s/\mathrm{Pa}) = 11.986 - \frac{1360\ \mathrm{K}}{T}
\]
\[
\lg(p_l/\mathrm{Pa}) = 9.729 - \frac{874\ \mathrm{K}}{T}
\]
计算(1) 二氧化碳三相点的温度和压力；
(2) 二氧化碳在三相点的熔化热和熔化熵。

\vspace{0.5cm}
\noindent\textbf{4. (15 pts.)}

硫酸铜的脱水反应$\ce{CuSO_{4}\cdot 3H_{2}O(s) = CuSO_{4}(s) + 3H_{2}O(g)}$在$25^\circ\mathrm{C}$时的标准平衡常数$K_p^\ominus = 10^{-6}$,为了使$0.01\ \mathrm{mol}\ \ce{CuSO_{4}}$转化为$\ce{CuSO_{4}\cdot 3H_{2}O(s)}$,问在$25^\circ\mathrm{C}$的$2\ \mathrm{dm}^3$容器中加入水蒸气的物质的量至少是多少。

\vspace{0.5cm}
\noindent\textbf{5. (15 pts.)}

某化合物的分解是一级反应,该反应活化能$E_\mathrm{a} = 14.43\times10^4\ \mathrm{J\cdot mol^{-1}}$。已知$557\ \mathrm{K}$时,该反应速率常数$k_1 = 3.3\times10^{-2}\ \mathrm{s^{-1}}$,现在要控制此反应在$10\ \mathrm{min}$内,转化率达到$90\%$,试问反应温度应控制在多少度。

\vspace{0.5cm}
\noindent\textbf{6. (15 pts.)}

反应$\ce{Zn(s) + CuSO_{4}(a=1) -> Cu(s) + ZnSO_{4}(a=1)}$在电池中进行。$15^\circ\mathrm{C}$时,测得$E=1.0934\ \mathrm{V}$,电池的温度系数$\left(\dfrac{\partial E}{\partial T}\right)_p = -4.29\times10^{-4}\ \mathrm{V\cdot K^{-1}}$,当电池有$2\ \mathrm{mol}$电子电荷量可逆输出时,
\begin{enumerate}[label=(\arabic*), itemsep=1ex]
    \item 写出电池表示式和电极反应式；
    \item 求电池反应的$\Delta_\mathrm{r}G_\mathrm{m}^\ominus$、$\Delta_\mathrm{r}S_\mathrm{m}^\ominus$、$\Delta_\mathrm{r}H_\mathrm{m}^\ominus$和$Q_r$；
    \item 若将该反应放在等压热量计内进行,求反应的热量$Q_p$。
\end{enumerate}